\section{Negative Results}
\label{sec:negative-results}

Not every idea we tested survived contact with held-out data. This
section documents six approaches that were implemented, evaluated
under the same gates as every shipped change
(Section~\ref{sec:validation}), and rejected. We report them in full
for two reasons. First, each rejection shaped the production system:
the discipline of adding one component at a time behind a held-out
gate is itself a consequence of the failures below. Second, several
of the hypotheses are natural enough that other practitioners will
try them; the measurements here may save that effort or sharpen it.
Table~\ref{tab:negative-summary} summarizes the six approaches; the
following subsections give one account each. Two failure modes recur:
constants derived from plausible first principles that measurement
contradicts, and signals that are genuinely informative but harmful
when injected at the wrong layer of the system.

\begin{table}[t]
  \centering
  \small
  \begin{tabular}{@{}p{0.26\linewidth}p{0.33\linewidth}p{0.33\linewidth}@{}}
    \toprule
    Approach & Hypothesis & Outcome \\
    \midrule
    Heuristic v2 component stack &
    Six domain-motivated bonuses improve the price-anchored heuristic &
    Trailed unmodified v1 in every backtested group round; reverted \\
    \addlinespace
    \texttt{team\_elo\_diff} as a GBM feature &
    The Elo gap adds fixture signal to the learned model &
    Held-out wash plus distribution-shift risk; signal routed to the
    structural backends instead \\
    \addlinespace
    Theoretical GK save-bonus multiplier (1.13) &
    A first-principles constant prices expected saves correctly &
    Overestimated by roughly a factor of two; empirical sweep
    shipped 0.50 \\
    \addlinespace
    \texttt{start\_rate\_lag} and \texttt{team\_gc\_form} as GBM
    features &
    Participation and team defensive form add predictive signal &
    Regressed leak-free walk-forward RMSE at every position; signals
    rehomed at the optimizer layer \\
    \addlinespace
    Benter combiner with market prices &
    Tournament-winner market probabilities complement model
    predictions &
    Any positive market weight raised RMSE for every backend in every
    round; shipped weight zero \\
    \addlinespace
    GK defensive-form adjustment &
    Recent goals conceded sharpens goalkeeper predictions &
    Worse as a feature and as a Poisson blend; structure beats
    realized form at this sample size \\
    \bottomrule
  \end{tabular}
  \caption{Summary of documented negative results. Each approach was
    implemented and evaluated against held-out or walk-forward data
    before rejection; details and evidence appear in the subsections
    below.}
  \label{tab:negative-summary}
\end{table}

\subsection{Heuristic v2: a six-component stack}
\label{sec:negative-heuristic-v2}

The heuristic backend of Section~\ref{sec:methodology} predicts
points as a position coefficient times price, scaled by matchup and
home factors: three moving parts. A v2 candidate layered six
domain-motivated changes onto it at once: a goalkeeper save bonus
scaled by opponent expected goals; a defender clean-sheet bonus
proportional to $\exp(-\mathrm{xG}_{\mathrm{opp}})$; ball-progression
and goal-conversion bonuses for midfielders and forwards who
outperform their position median; a premium-tier bump for players
priced above \$9.5 million; and a realized-form anchor that blends
the structural prediction toward the player's tournament scoring
average, with a weight growing from zero to 0.40 by four matches
played. Every component was individually defensible, and several
encode genuine features of the scoring rules.

Replayed under the frozen-data backtest protocol of
Section~\ref{sec:live-results}, v2 scored below unmodified v1 in
every completed group round (Table~\ref{tab:backtest-rounds}). Our
diagnosis has four parts. First, the realized-form anchor pulls
toward one to three matches of data, where the per-player average is
dominated by variance; the structural prediction is more reliable
than the realized average at that sample size. Second, the stacked
bonuses shifted squad composition, promoting players (for example,
defenders facing low-xG opponents) who did not in fact outscore the
v1 picks. Third, the premium-tier bump was too small to change the
optimizer's selections and should have been either amplified or
omitted. Fourth, and most generally, six additions give the model six
places to be slightly wrong where v1 has three; on small samples the
misalignments compound rather than cancel.

We reverted to v1 and kept v2 in the repository as a runnable
research backend. We note for completeness that the replay ordering
reversed in later knockout rounds (Table~\ref{tab:backtest-rounds}),
with cumulative totals of \TotalHeuristicVTwo{} for v2 against
\TotalHeuristic{} for v1, which underlines how noisy a three-round
gate is. \POSTFINAL{state whether the v2 knockout-round recovery
persisted over the full tournament and reconcile with the reversion
decision} The procedural lesson stood regardless and became policy:
one component at a time, each behind a held-out gate, with failures
documented rather than discarded.

\subsection{\texttt{team\_elo\_diff} as a GBM feature}
\label{sec:negative-elo-feature}

The gradient-boosted model \citep{ke2017lightgbm} shipped with five
features. A v3 candidate added the Elo gap between a player's team
and the opponent \citep{elo1978rating, hvattum2010using} as a sixth:
club Elo computed from historical match data on the training side,
country Elo aliased to the same column at inference. Under a
deterministic seed, the held-out comparison against v2 was a wash:
small improvements for goalkeepers and midfielders, offset by small
regressions for defenders and forwards, with all four movements
within about one percent.

A wash on its own might still justify shipping, but the feature
carried a structural risk that the aggregate number understates:
distribution shift. Elo gaps between national teams at a World Cup
are several times wider than the club-Elo gaps in the training
league, so the tree ensemble would evaluate many tournament fixtures
in a region of feature space with no training support, where trees
extrapolate by clamping to the nearest leaf. We did not ship v3. The
Elo signal is instead consumed directly by the heuristic and Poisson
backends, where its effect on predictions is monotone by construction
and no extrapolation is involved. A procedural footnote: this A/B was
only legible after we fixed all training seeds, since run-to-run
noise from the boosting procedure's subsampling initially exceeded
the effect being measured.

\subsection{The theoretical goalkeeper save-bonus multiplier}
\label{sec:negative-gk-multiplier}

The Poisson backend's original goalkeeper formula assigned a flat
expected save bonus per start, independent of the fixture. Replacing
the constant with a term scaled by opponent expected goals is clearly
right in direction: the save-bonus rule pays one point per three
saves, and save opportunities are a function of how much attack the
opposing team generates. The question is the multiplier. We first
derived it from football fundamentals: roughly 4.0 shots on target
per unit of opponent xG, a median save percentage of 0.85, and the
one-per-three-saves rule give $4.0 \times 0.85 / 3 \approx 1.13$
expected bonus points per unit of opponent xG.

The derivation was internally consistent and built from three
independently measurable quantities. It was also wrong. On held-out
data the 1.13 multiplier regressed goalkeeper RMSE relative to the
flat baseline it was meant to replace
(Figure~\ref{fig:gk-sweep}). Both inputs were overestimates: much of
opponent xG never becomes an on-target shot, and the median
goalkeeper in the evaluation distribution saves at a substantially
lower rate than the elite keepers who anchor intuition. The two
errors compound multiplicatively. An empirical sweep over candidate
multipliers found its minimum at 0.50, less than half the derived
value, improving goalkeeper RMSE from \GkRmseVOne{} to \GkRmseVTwo{}
(Table~\ref{tab:gk-ab}); the empirical calibration and its A/B
validation are discussed in Section~\ref{sec:analysis}. The lesson we
draw is not that structural derivation is useless (the functional
form survived; only the constant fell) but that a derivation is a
hypothesis, and the sweep, not the algebra, sets the constant.

\subsection{Participation and defensive form as GBM features}
\label{sec:negative-minutes}

A benched player scores at most the appearance floor regardless of
form or matchup, so participation is an obvious candidate signal.
With no reliable predicted-lineup source for the tournament, we
proxied it: \texttt{start\_rate\_lag} is the trailing fraction of a
player's recent completed rounds with a positive score (the true
starts flag is used on the league training side), computed leak-free
with a one-round shift. The signal is real: players with full recent
participation outscore rotation risks by a wide margin and blank far
less often. \texttt{team\_gc\_form}, the trailing mean of goals
conceded by the player's team, was constructed analogously as a
defensive-form signal.

Feeding both to the GBM failed. Configuration D of the walk-forward
validation added the pair to the shipped feature set and regressed
the leak-free pooled RMSE from \WalkforwardPooledC{} to
\WalkforwardPooledD{}, worse at every position
(Table~\ref{tab:walkforward}). Over a handful of tournament rounds
the participation proxy is too noisy, and it is collinear with the
trailing form feature (a benched player already shows low form), so
it adds variance the trees cannot convert into signal; defensive form
is likewise collinear with the strength features the model already
holds. The correct layer turned out to be the optimizer, not the
model: an availability discount scales a player's effective points by
$0.5 + 0.5 \cdot \mathtt{start\_rate\_lag}$ before squad selection,
so the solver stops spending budget on rotation risks. That
adjustment improved realized squad points over the backtested rounds
(\POSTFINAL{number: realized squad points with versus without the
availability discount on the heuristic backend}) and was stable
across a wide band of floor values; it also surfaces a rotation-risk
flag in the operator report. The same walk-forward harness rejected
three further feature-layer candidates in the same review: a
form-times-opponent-strength interaction (within noise), recency
weighting of tournament training rows (monotonically worse for
goalkeepers), and dropping did-not-play rows to match the league
pipeline's minutes filter, which was decisively worse because the
zero-target rows are precisely how the model learns that
low-participation profiles score zero.

\subsection{A Benter-style combiner over prediction-market prices}
\label{sec:negative-market}

\citet{benter1994computer} showed that combining a fundamental model
with market odds can beat either alone, and prediction markets are
efficient aggregators of public information
\citep{wolfers2004prediction}. We therefore tested a combiner that
multiplies each player's model prediction by a country-level
adjustment derived from tournament-winner contract prices on a
prediction market: the most recent pre-deadline implied win
probability per country, normalized by the field mean, entering
through a single free weight $\beta_2$ (intercept and slope fixed at
zero and one). We swept ten values of $\beta_2$ from 0 to 0.30 for
every backend over the completed rounds.

The result was unambiguous. The empirical optimum for every backend
is $\beta_2 = 0$, and RMSE rises monotonically in the market weight.
At the ablation operating point $\beta_2 = \MarketBetaTwo$, enabling
the adjustment worsens RMSE for every backend in every round
(Table~\ref{tab:market-negative}), by \MarketMeanDeltaRmse{} on
average across backends and rounds; Figure~\ref{fig:market-negative}
shows the with-versus-without comparison.

\begin{figure}[t]
  \centering
  \includegraphics[width=0.85\textwidth]{fig_market_negative.pdf}
  \caption{Effect of blending prediction-market signals into the
    expected-points models: RMSE with versus without the market term.
    \figsource{scripts/with\_vs\_without\_market.py $\rightarrow$
    data/evaluation/with\_vs\_without\_market.json}}
  \label{fig:market-negative}
\end{figure}

\begin{table}[t]
  \centering
  % GENERATED by scripts/paper_tables.py -- do not edit
\begin{tabular}{lrrrr}
\toprule
Backend & R1 & R2 & R3 & Mean \\
\midrule
Heuristic & 0.364 & 0.338 & 0.266 & 0.323 \\
Poisson & 0.938 & 1.111 & 1.084 & 1.044 \\
GBM & 0.201 & 0.163 & 0.150 & 0.171 \\
Monte Carlo & 0.866 & 1.730 & 1.230 & 1.275 \\
\bottomrule
\end{tabular}

  \caption{Per-round change in RMSE ($\Delta$RMSE, positive is worse)
    when the Benter-style market adjustment is enabled
    ($\beta_2 = \MarketBetaTwo$).
    \figsource{scripts/with\_vs\_without\_market.py}}
  \label{tab:market-negative}
\end{table}

We attribute the failure primarily to a granularity mismatch.
Tournament-winner contracts price a season-scale outcome driven by
squad depth, bracket, and long-horizon strength; our target is
per-match, per-player fantasy points, driven by fixture-scale
factors the contract cannot see. What country-level strength the
market does encode is redundant with signals the models already
carry: the implied win probabilities correlate at
\POSTFINAL{number: correlation between market implied win
probabilities and derived country Elo} with our derived country Elo,
making the market term a noisy near-substitute rather than a
complement. Two further effects plausibly contribute: the prices
exhibit the canonical favorite-longshot bias
\citep{thaler1988favorite, snowberg2010explaining}, which our
multiplicative combiner does not correct, and liquidity is
concentrated on a few favorites, so the thin long-shot contracts we
weight symmetrically are dominated by noise. The negative result is
specific to this configuration: tournament-winner contracts at
country granularity, a per-player per-match target, and a
multiplicative combiner. Per-fixture markets at granularities matched
to the prediction target remain untested and are discussed in
Section~\ref{sec:future_work}. We kept the combiner architecture and
the market scraper, shipped a weight of zero (the combiner is a
no-op in production), and demoted market data to descriptive context
on the intelligence page.

\subsection{Goalkeeper defensive-form adjustment}
\label{sec:negative-gk-form}

The last negative result closed the loop on goalkeeper form. Our
working hypothesis was that personal scoring form is noise for
goalkeepers (their points come from clean sheets and saves, not from
repeatable individual scoring) and that the right form signal for the
position is the team's recent defensive record. The first half held
up; the second did not. We tested \texttt{team\_gc\_form} in two
ways: as a GBM feature, where it participated in the configuration D
regression of Section~\ref{sec:negative-minutes}, and blended
directly into the Poisson clean-sheet term, where goalkeeper points
actually originate. Sweeping the blend weight on held-out goalkeeper
rows, RMSE rose monotonically with the weight: the more realized
defensive form we injected, the worse the prediction.

The mechanism is the same one that sinks personal form for the
position. Three or four matches of goals conceded is a tiny,
high-variance sample, and its dominant component is opponent
strength, which the structural model already encodes through the
fixture terms \citep{maher1982modelling, dixon1997modelling}. The
structural Poisson backend remains our best goalkeeper predictor
(Section~\ref{sec:validation}), the ensemble routes goalkeepers to
it, and no defensive-form signal ships. The honest finding is that at
tournament sample sizes no realized-form signal we tested beats
structure for this position, and pretending otherwise would have cost
accuracy. Section~\ref{sec:lessons} generalizes the pattern.
